\documentclass[11pt,a4paper]{article}
\usepackage[utf8]{inputenc}
\usepackage[T1]{fontenc}
\usepackage{amsmath,amssymb}
\usepackage{graphicx}
\usepackage{booktabs}
\usepackage{authblk}
\usepackage{hyperref}
\usepackage{geometry}
\geometry{margin=1in}
\usepackage{siunitx}
\usepackage{xcolor}

\title{BELA: A Fully Automated Deep Learning Framework for Ploidy Prediction and
Quality Assessment of Human Blastocysts Using Time-Lapse Imaging}

\author[1]{Anonymous Author}
\affil[1]{Department of Computational Biology}

\date{}

\begin{document}
\maketitle

\begin{abstract}
Assessing fertilized human embryos is a central but labor-intensive task in
in-vitro fertilization (IVF) and remains the dominant bottleneck in
selecting embryos with the highest implantation potential. Existing
artificial intelligence models for embryo quality assessment and ploidy
prediction do not fully exploit the rich temporal information contained in
time-lapse imaging, and most still require subjective input from
embryologists. We present BELA (Blastocyst Evaluation Learning Algorithm),
a fully automated, multitask deep learning pipeline that predicts embryo
ploidy status directly from day-5 time-lapse video sequences and maternal
age, without any manual annotation. BELA first uses a multitask
bidirectional long short-term memory (BiLSTM) network to regress the
embryologist-derived blastocyst score together with its inner cell mass
(ICM), trophectoderm (TE), and expansion sub-components from VGG16
features extracted between 96 and 112 hours post insemination (hpi). The
model-derived blastocyst score (MDBS) is then passed, along with maternal
age, to a logistic regression classifier that discriminates euploid (EUP)
from aneuploid (ANU) or complex aneuploid (CxA) embryos. BELA achieves a
mean absolute error of $1.855 \pm 0.03$ between the MDBS and the ground
truth blastocyst score on the WCM-Embryoscope test set, and attains an
area under the receiver operating characteristic curve (AUC) of $0.66 \pm
0.008$ for EUP vs.\ ANU discrimination, which rises to $0.76 \pm 0.002$
upon inclusion of maternal age. For the EUP vs.\ CxA task, BELA reaches an
AUC of $0.708 \pm 0.004$, increasing to $0.826 \pm 0.004$ when maternal
age is added. BELA is competitive with a classifier trained on
embryologist-annotated blastocyst scores and significantly outperforms a
day-5 video baseline across all test sets ($p < 0.05$). While BELA is not
intended to replace preimplantation genetic testing for aneuploidy
(PGT-A), it provides a standardized, non-invasive, low-cost decision
support tool that can streamline embryo selection in IVF clinics. BELA is
deployed as a web application, STORK-V, that takes time-lapse frames
between 96 and 112 hpi and maternal age as input and returns calibrated
ploidy probabilities together with interpretable intermediate quality
scores.
\end{abstract}

\section{Introduction}\label{sec:intro}

Since the advent of in-vitro fertilization (IVF) in 1978, assisted
reproductive technology has supported the birth of over 8 million babies
worldwide and has become a cornerstone of infertility treatment
\cite{dyer2016international}. The procedure involves transvaginal transfer
of laboratory-fertilized oocytes into the uterus, and its success depends
critically on the selection of a single high-quality embryo with normal
chromosomal content \cite{practice2017performing,munne2019preimplantation}.
Embryos with the correct chromosomal constitution (\emph{euploid}) are
associated with higher implantation rates and successful pregnancies,
whereas \emph{aneuploid} embryos are linked with miscarriage, failed
implantation, and chromosomal disorders such as Down or Turner syndrome.
Embryo aneuploidy, and with it the risk of miscarriage, increases sharply
with advanced maternal age \cite{franasiak2014nature}.

The clinical gold standard for ascertaining embryo ploidy is
preimplantation genetic testing for aneuploidy (PGT-A), which requires a
biopsy of trophectoderm (TE) cells, whole genome amplification, and
chromosomal copy number analysis. Despite its ability to improve
implantation rates, PGT-A is invasive, costly, and time-consuming, and can
compromise embryo viability \cite{munne2019preimplantation}. Its accuracy
is further limited by embryonic mosaicism, where aneuploid and euploid
cells coexist in the TE, producing false calls and reduced implantation
rates \cite{popovic2020mosaic}.

In parallel, the accumulation of large time-lapse imaging datasets and the
maturation of computer vision methods have enabled fully automated embryo
quality assessment. Khosravi et al.\ introduced STORK, a deep learning
model that predicts embryo morphology and correlates with successful birth
outcomes \cite{khosravi2019deep}. More recently, several groups have
attempted to repurpose similar architectures for ploidy prediction.
Chavez-Badiola et al.\ developed ERICA to analyze 1{,}231 embryo images
for ploidy prediction, achieving 70\% accuracy, an AUC of 74\%, a
sensitivity of 54\%, and a specificity of 86\%, and correctly identifying
a euploid embryo as top-ranked in 79\% of cases
\cite{chavezbadiola2020predicting}. Nevertheless, ERICA could not
distinguish between single and complex aneuploidies and relied on a
limited dataset. Barnes et al.\ trained machine learning algorithms to
predict ploidy from a single image at 110 hours post insemination (hpi)
using time-lapse sequences \cite{barnes2023noninvasive}, while the UBar
CNN--LSTM model of Silver et al.\ attempted to exploit full time-lapse
sequences and reached an AUC of 0.82 on a limited dataset
\cite{silver2020data}. Lee et al.\ trained a two-stream inflated 3D model
on 670 image sequences and reported an AUC of 0.74 for discriminating
euploid/mosaic from aneuploid embryos \cite{lee2022end}.

A persistent challenge in ploidy prediction from time-lapse imaging is
that not all developmental stages carry the same amount of information
about chromosomal status. Many prior works therefore restrict attention to
hand-picked developmental windows
\cite{campbell2013modelling,kragh2019automatic}. For example, Campbell et
al.\ proposed blastocyst expansion timing on day 5 as a predictor of
ploidy \cite{campbell2013modelling}, but the predictive accuracy of such
morphokinetic rules varies substantially across clinics
\cite{kramer2014assessing}. Analyzing full embryo development videos would
in principle circumvent the need for hand-tuned time windows, but the
computational cost of training on vast raw video datasets and the noise
contributed by irrelevant frames can quickly compromise performance.

We address these limitations by introducing BELA, the Blastocyst
Evaluation Learning Algorithm---a fully automated model that predicts
ploidy status from time-lapse sequences and maternal age alone. BELA
contributes: (i) a multitask BiLSTM that jointly regresses the blastocyst
score and its ICM, TE, and expansion components from VGG16 features
extracted from day-5 time-lapse frames, yielding a
\emph{model-derived blastocyst score} (MDBS) that correlates with the
embryologist-derived score and, on external data, with ploidy; (ii) a
logistic regression classifier that combines MDBS with maternal age to
predict EUP vs.\ ANU and EUP vs.\ CxA ploidy classes; (iii) a systematic
comparison with a direct day-5 video baseline and with classifiers trained
on embryologist-derived blastocyst scores, on four datasets from Weill
Cornell Medicine, IVI Valencia, and IVF Florida; and (iv) STORK-V, a
web-based clinical tool that packages BELA into a reproducible deployment
for embryologists. By removing the need for subjective manual annotation
and by focusing on the informative 96--112 hpi window, BELA streamlines
ploidy prediction and is broadly applicable across different clinical
settings.

\section{Methods}\label{sec:methods}

\subsection{Datasets}

We used four time-lapse datasets to develop and validate BELA. The
\emph{WCM-Embryoscope} dataset, collected at the Center for Reproductive
Medicine at Weill Cornell Medicine (WCM) between 2018 and 2019, comprises
1{,}998 time-lapse sequences from the Embryoscope\textsuperscript{\textregistered}
imaging system, with PGT-A-derived labels distributed as 916 euploid (EUP),
494 single aneuploid (SA), and 588 complex aneuploid (CxA) embryos. A
total of 498 patients are represented, with an average of four biopsied
embryos per patient; each sample was treated independently. The
\emph{WCM-Embryoscope+} dataset, collected between 2019 and 2020 using the
newer Embryoscope+\textsuperscript{\textregistered} system, comprises 841
sequences (410 EUP, 170 SA, 261 CxA). Each WCM dataset provides the
embryologist-derived blastocyst score (BS), its ICM, TE, and expansion
sub-components, morphokinetic parameters, and maternal age at oocyte
retrieval. The blastocyst score is formulated as in prior WCM work and
takes integer values between 3 and 14, with lower scores indicating higher
quality \cite{kragh2019automatic}. Two external datasets were used to
assess generalization: the \emph{Spain} dataset from IVI Valencia (543
embryos: 234 EUP and 309 ANU) and the \emph{Florida} dataset from IVF
Florida (869 embryos: 445 EUP, 202 SA, 222 CxA). Spain data provide only
EUP/ANU labels, while Florida data additionally distinguish SA and CxA
embryos and include maternal age and blastocyst score. Dataset
characteristics are summarized in Table~\ref{tab:datasets}.

\begin{table}[t]
\centering
\caption{Characteristics of training and validation datasets. Ploidy
classes: euploid (EUP), single aneuploid (SA), complex aneuploid (CxA),
aneuploid (ANU, Spain only combines SA and CxA).}
\label{tab:datasets}
\begin{tabular}{lrrrrl}
\toprule
Dataset & $N$ & EUP & SA & CxA & Imaging instrument \\
\midrule
WCM-Embryoscope  & 1{,}998 & 916 & 494 & 588 & Embryoscope\textsuperscript{\textregistered} \\
WCM-Embryoscope+ &    841 & 410 & 170 & 261 & Embryoscope+\textsuperscript{\textregistered} \\
Spain (IVI Valencia) & 543 & 234 & \multicolumn{2}{c}{309 (ANU)} & Embryoscope\textsuperscript{\textregistered} \\
Florida (IVF Florida) & 869 & 445 & 202 & 222 & Embryoscope+\textsuperscript{\textregistered} \\
\bottomrule
\end{tabular}
\end{table}

\subsection{Preimplantation Genetic Testing}

Embryos from Weill Cornell were biopsied on day 5 or day 6 depending on
when they reached blastocyst stage, and biopsied cells were analyzed by
next generation sequencing (NGS) at the Ronald O.\ Perelman and Claudia
Cohen Center for Reproductive Medicine. For the Spain dataset, Igenomix
Spain performed PGT-A; embryos were subjected to assisted hatching on
day 3 with a Hamilton-Thorne Lykos\textsuperscript{\textregistered} laser,
and 5--6 trophectodermal cells were biopsied after blastocyst formation
and assessed with Thermo Fisher Scientific NGS technology. Florida samples
were analyzed using the same Igenomix/Thermo Fisher NGS workflow
\cite{garcia2020optimized}.

\subsection{Temporal and Spatial Processing of Time-Lapse Sequences}

Raw time-lapse sequences varied considerably in length, frame rate, start
and end points, and missing frames, which could bias video classifiers if
left uncorrected. We therefore standardized all sequences on a uniform
time grid. Standardized time points were defined at 30-minute intervals
from 0 to 150 hpi (i.e., 0 hpi, 0.5 hpi, \ldots, 149.5 hpi, 150.0 hpi),
yielding 301 frames per embryo. For each embryo, the time-lapse image
closest to each standardized time point was assigned to that point; if no
image fell within a 2-hour window of a standardized time point, a blank
frame was assigned. A 2-hour window was chosen because (i) significant
embryonic changes typically occur at intervals greater than 2 hours, so
this window captures important dynamics while still allowing
standardization, and (ii) it matches the typical frame acquisition rate of
the imaging systems used. Once standardized sequences were constructed,
frames could be extracted for downstream modelling using three parameters:
\texttt{start\_hr}, \texttt{end\_hr}, and \texttt{interval}. For example,
a model trained on day-2 embryo development would use
\texttt{start\_hr}$=24.0$ hpi, \texttt{end\_hr}$=48.0$ hpi and
\texttt{interval}$=2$ hrs, which yields 13 frames. For image
classification experiments, a single time point could be selected and the
corresponding frame extracted.

Sequences rendered unusable by standardization were excluded from
analysis. Exclusion criteria included cases where the embryo was absent
from the petri dish, less than half-visible, or the image was too dim to
discern the embryo. We resized each frame from $800 \times 800$ to
$224 \times 224$ pixels. To reduce background bias, we applied a circle
Hough Transform for embryo segmentation in each frame. This pipeline was
applied uniformly to the WCM-Embryoscope, WCM-Embryoscope+, Spain, and
Florida datasets. To increase the diversity and robustness of the
training data we applied video augmentation, including random horizontal
flipping (which doubles the effective dataset size by producing mirror
images) and random rotations (which simulate varied embryo orientations).

\subsection{Feature Extraction}

Spatial features were extracted from each frame with an ImageNet
pretrained VGG16 convolutional neural network from TensorFlow 2.7
\cite{simonyan2014very}. The final average pooling layer of VGG16 yields a
512-dimensional feature vector per frame, so that each embryo is
represented by a $T \times 512$ matrix where $T$ is the number of frames
retained from its standardized sequence. We explored alternative
pretrained feature extractors but found that the ImageNet-pretrained VGG16
consistently provided the best downstream performance.

\subsection{BELA Prediction Pipeline}

BELA is organized as a two-step pipeline. In
the first step, a multitask bidirectional LSTM (BiLSTM) regression model
takes the 512-dimensional per-frame features from the day-5 time window
(96--112 hpi) and jointly regresses the overall blastocyst score (BS), the
ICM score, the TE score, and the expansion score. The architecture
consists of one bidirectional LSTM layer followed by two multi-unit dense
layers; each prediction head is a single-unit dense layer. Because all
four targets are continuous, we use the log-cosh loss and the Adam
optimizer, with equal loss weights across tasks. Although ICM, TE, and
expansion score are the components of BS, multitask training allows the
BiLSTM to exploit overlapping representations and improves BS prediction.
We also explored attention mechanisms and stacking multiple bidirectional
LSTM layers, but neither improved performance significantly ($p > 0.05$).
Maternal age was included as an auxiliary input feature in the BiLSTM.
Early stopping with patience$=5$ was used to prevent overfitting by
monitoring the validation loss.

In the second step, the predicted blastocyst score from the BiLSTM---the
\emph{model-derived blastocyst score} (MDBS)---is fed, together with
maternal age, into a logistic regression classifier to predict ploidy
status. The logistic regression used cross-entropy loss. Two prediction
tasks were modelled: euploid vs.\ aneuploid (EUP vs.\ ANU) and euploid vs.\
complex aneuploid (EUP vs.\ CxA).

\subsection{Training and Evaluation Protocol}

The BiLSTM regression model was trained only on the training slice of the
WCM-Embryoscope dataset. To prevent data leakage, the WCM-Embryoscope
dataset was split 70/30 into training and test partitions. Exclusion
criteria reduced the dataset from 1{,}998 to 1{,}684 embryos. Four-fold
cross-validation was applied on the training split; for each fold, a
BiLSTM regression model was trained and a corresponding logistic
regression was fit on the fold's MDBS and maternal age. Performance
metrics from the four logistic regression models were averaged, and
standard deviations across folds are reported throughout. The first
component of BELA was evaluated using the mean absolute error (MAE) of
the MDBS relative to the embryologist-derived BS; the second component
was evaluated using accuracy, AUC, precision, and recall. For comparison,
we trained two baseline models on the same cross-validation splits: a
\emph{day-5 video model} that directly predicts ploidy from time-lapse
input in the 96--112 hpi window using a BiLSTM, and an
\emph{embryologist-annotated model} that uses only the ground-truth
blastocyst score in a logistic regression.

All statistical comparisons between model means were made using the
Student's $t$-test, which is appropriate for approximately normal data
with similar variances. Because many comparisons were performed, we
applied a Bonferroni correction to control the family-wise error rate.

\subsection{Computational Resources}

Model training and inference were performed on an Apple M1 Mac with
TensorFlow Metal. Logistic regression models trained in $2.5 \pm 1.2$
seconds on average, whereas BiLSTM models required $30.3 \pm 11$ minutes.
For deployment, the BELA model used by the STORK-V web application was
trained on a high-performance BioHPC computing cluster at Cornell, Ithaca,
using an NVIDIA A40 GPU; training completed in $5.23$ minutes. Inference
for a single embryo on the STORK-V platform took $30 \pm 5$ seconds.

\section{Results}\label{sec:results}

\subsection{Multitask Blastocyst Score Regression}

The first component of BELA was trained to regress the blastocyst score
and its ICM, TE, and expansion sub-components from day-5 time-lapse
features. Across both training and test splits of the WCM-Embryoscope
dataset, the MDBS produced by the multitask BiLSTM exhibited a moderate
correlation with the embryologist-derived BS (Pearson correlation
$>0.7$). Similar moderate correlations were observed between predicted
and actual ICM, TE, and expansion scores. The mean absolute error between
the MDBS and the ground-truth WCM-Embryoscope blastocyst scores was
$1.855 \pm 0.03$. Replacing the multitask BiLSTM with a single-task BiLSTM
that regressed only BS increased the MAE to $1.877 \pm 0.027$ on the same
test set, supporting the benefit of joint training across all four
blastocyst score targets.

\subsection{Ploidy Classification on Weill Cornell Data}

BELA and its two baselines were evaluated on the WCM-Embryoscope test set
for the EUP vs.\ ANU and EUP vs.\ CxA tasks, both with and without
maternal age as an additional input (Table~\ref{tab:bela_auc}). When
trained to distinguish EUP from ANU, BELA attained an AUC of $0.66 \pm
0.008$, which rose to $0.76 \pm 0.002$ upon inclusion of maternal age.
For the EUP vs.\ CxA task, BELA reached an AUC of $0.708 \pm 0.004$
without age and $0.826 \pm 0.004$ with maternal age. In all tested
scenarios---including or excluding age, test sets, and prediction
tasks---BELA outperformed the day-5 video baseline that directly
classifies ploidy from 96--112 hpi frames ($p < 0.05$).

\begin{table}[t]
\centering
\caption{BELA performance on the WCM-Embryoscope test set. AUC values are
reported as mean $\pm$ standard deviation across four cross-validation
folds.}
\label{tab:bela_auc}
\begin{tabular}{lll}
\toprule
Task & Input & AUC \\
\midrule
EUP vs.\ ANU & MDBS only & $0.66 \pm 0.008$ \\
EUP vs.\ ANU & MDBS + maternal age & $0.76 \pm 0.002$ \\
EUP vs.\ CxA & MDBS only & $0.708 \pm 0.004$ \\
EUP vs.\ CxA & MDBS + maternal age & $0.826 \pm 0.004$ \\
\bottomrule
\end{tabular}
\end{table}

Without maternal age, the embryologist-annotated blastocyst score baseline
surpassed BELA in all prediction tasks except EUP vs.\ ANU on the
WCM-Embryoscope test set ($p < 0.05$), reflecting the fact that the
embryologist-annotated BS is itself a highly calibrated summary of day-5
morphology. With maternal age incorporated, however, BELA significantly
outperformed the embryologist-annotated BS model on the WCM-Embryoscope
test set ($p < 0.05$), while still modestly underperforming on the
WCM-Embryoscope+ dataset.

\subsection{External Validation on the Spain Dataset}

We next applied BELA to the Spain dataset (IVI Valencia, 543 embryos),
which contains only EUP/ANU labels. Embryos in the Spain dataset
underwent assisted hatching on day 3, which affected later blastocyst
morphology and morphokinetics: zona pellucida was frequently bleached and
full blastocyst expansion was not reached. Despite these systematic
differences, BELA significantly outperformed the day-5 video baseline
both with and without maternal age ($p < 0.05$). To quantify the
domain shift, we averaged VGG16 per-frame features across the 96--112 hpi
window for each embryo and reduced them with PCA. Embryos from the US
clinics clustered together while Spanish embryos clustered distinctly in
the lower-right of the feature space, yet BELA's AUC without maternal
age remained comparable to that obtained on Weill Cornell data. This
observation suggests that BELA's reliance on the MDBS as an intermediate
representation confers a degree of domain robustness relative to a
video-only baseline. The models that also incorporated maternal age,
however, showed reduced performance on the Spain dataset, most likely
because of substantial differences in the maternal age distribution
between Spain and Weill Cornell cohorts.

\subsection{External Validation on the Florida Dataset}

On the Florida dataset (869 embryos: 445 EUP, 202 SA, 222 CxA), the
absolute performance of BELA and of the comparison models (logistic
regression on maternal age and/or embryologist-derived BS) decreased
relative to the WCM test sets. Two factors likely contributed to this
decrease. First, in the Florida cohort, both maternal age and the
embryologist-derived blastocyst score were only weakly correlated with
ploidy, reducing the useful signal in the strongest predictors of all
models. Second, Florida blastocyst scores were predominantly centered
around a value of 7, reducing their granularity as a predictor. Florida
clinicians evaluate blastocysts at 115 and 144 hpi, in contrast with the
Weill Cornell protocol which also considers earlier timepoints, and this
difference in scoring practice likely drives the lack of granularity.

Interestingly, the MDBS produced by the first module of BELA exhibited a
slightly stronger correlation with ploidy status on the Florida dataset
(Pearson correlation $-0.119$) than the embryologist-derived blastocyst
score ($-0.101$), suggesting that BELA's learned score better captures
ploidy-relevant signal than the hand-labelled score on this dataset. To
validate this observation, a Weill Cornell embryologist re-graded the 50
Florida embryos for which the MDBS deviated most from the Florida
blastocyst scores. The mean absolute error between the MDBS and the
re-graded Weill Cornell blastocyst scores was $4.16$, compared with
$5.02$ between the MDBS and the original Florida scores, indicating a
higher agreement between MDBS and the more granular Weill Cornell
scoring scheme. Consistent with this, BELA (with maternal age)
significantly outperformed a logistic regression trained on maternal age
and the embryologist-derived Florida blastocyst score for the EUP vs.\
ANU task ($p < 0.05$).

\subsection{STORK-V Web Deployment}

To make BELA available to embryologists in clinical settings, we developed
STORK-V, a web-based application that accepts time-lapse images captured
between 96 and 112 hpi and the maternal age as inputs, and returns
calibrated probabilities of euploidy, aneuploidy, and complex aneuploidy,
together with intermediate quality scores that can support further
inspection by the embryologist. Two models are deployed behind the
interface: one trained for the EUP vs.\ ANU task and one for the EUP
vs.\ CxA task. STORK-V is intended as a decision-support tool that gives
clinicians a convenient, consistent, and non-invasive way to assess
embryo ploidy status as part of IVF workflows.

\section{Discussion}\label{sec:discussion}

This work introduces BELA, a fully automated, multitask deep learning
framework that predicts blastocyst quality scores and embryo ploidy from
day-5 time-lapse imaging and maternal age alone. Unlike prior work that
uses either still images or direct video classification, BELA exploits
the intermediate embryologist-derived blastocyst score as a soft
supervision target: the first component of BELA jointly regresses BS,
ICM, TE and expansion scores from VGG16 features of frames in the
96--112 hpi window, and the second component turns the resulting MDBS
into ploidy probabilities via a simple logistic regression that also
uses maternal age. This design yields several practical advantages.
First, BELA is competitive with a classifier trained on
embryologist-annotated blastocyst scores and significantly outperforms a
direct day-5 video model on all test sets. Second, by focusing on the
96--112 hpi window, BELA avoids the computational and statistical cost
of modelling entire multi-day sequences. Third, because BELA consumes
only time-lapse frames and maternal age, it can be dropped into existing
IVF workflows without disrupting established protocols. Fourth, BELA
offers a degree of explainability: embryologists can inspect the MDBS
and the predicted ICM, TE, and expansion sub-scores to understand why a
given embryo was classified as euploid or aneuploid.

Across SART age groups, BELA's AUC is bimodal on the WCM-Embryoscope and
WCM-Embryoscope+ datasets, with the best performance at the lower and
higher age extremes. On the Spain and Florida datasets, the
age-stratified performance does not follow the WCM pattern, reflecting
demographic differences across clinics: IVF is more affordable in Spain,
which broadens the population seeking treatment, while in the United
States the high cost of IVF restricts access. Interestingly, BELA's
EUP vs.\ CxA model tends to split single aneuploid (SA) embryos roughly
evenly between EUP and CxA predictions, suggesting that SA embryos often
resemble either extreme and are intrinsically harder to classify. This
observation is consistent with BELA models that are explicitly trained
to discriminate EUP from SA. BELA is competitive even though its test
sets consist exclusively of high-quality biopsied embryos, a harder
benchmark than the common practice of including non-viable embryos as
negatives which tends to inflate AUCs.

The study has several limitations. Video classification models are data
hungry, and with only $\sim2{,}000$ training sequences we were unable to
experiment with more expressive architectures such as 3D ConvNets or
two-stream inflated networks. Although we tried several pretrained
feature extractors, none outperformed the ImageNet-pretrained VGG16; a
more appropriate backbone might reveal signal from earlier developmental
stages. We also lacked access to potentially informative maternal
features such as hormone levels at oogenesis, demographics, and other
clinical variables. BELA is trained to predict the blastocyst score as
an intermediate label, and although this score is well-documented and
predictive of ploidy, it remains manually curated and subject to
intra-observer variability. BELA does not explicitly model mosaic
embryos, which we deliberately excluded; as a result mosaic embryos with
high implantation potential could be misclassified. Finally, differences
in PGT-A platforms and biopsy protocols across clinics \cite{capalbo2022}
introduce variability in the ground-truth labels that can limit
generalizability.

Despite these caveats, BELA provides a standardized, non-invasive,
cost-effective framework for embryo selection and prioritization. It is
not intended to replace PGT-A, but rather to support it: by identifying
embryos with a high probability of being euploid, BELA can reduce the
number of unnecessary biopsies, streamline embryologist workflows, and
lower patient costs. By flagging embryos likely to be complex aneuploid,
BELA can help clinicians focus attention on cases where embryo arrest or
implantation failure is most likely, which is particularly valuable for
patients with few available embryos. Automated blastocyst score
prediction is also clinically useful on its own, as it can standardize
blastocyst scoring across clinics and remove embryologist subjectivity.
Future iterations of BELA should explore richer backbones, larger
multi-center training corpora, explicit modelling of mosaic embryos,
integration of additional maternal covariates, and prospective clinical
validation of STORK-V.

\bibliographystyle{unsrt}
\bibliography{references}

\end{document}
